\begin{proposition}[完全正仪器、效应算符与条件态]
在有限工作空间中，若测量结果 $x$ 的更新由 Kraus 算符
$M_{x\alpha}$ 给出，且 $\sum_{x,\alpha}M_{x\alpha}^\dagger M_{x\alpha}=I$，则
\[
\mathcal I_x(\rho)=\sum_\alpha M_{x\alpha}\rho M_{x\alpha}^\dagger,
\qquad \Pi_x=\sum_\alpha M_{x\alpha}^\dagger M_{x\alpha}
\]
构成量子仪器和 POVM。每个分支完全正且迹不增，
$p_x=\Tr(\Pi_x\rho)$，所有分支之和保迹。
\end{proposition}
\begin{proof}
将任意参考系统 $R$ 加入，对任意联合正态 $X_{SR}$，
\[
(\mathcal I_x\otimes\mathcal I_R)(X_{SR})
=\sum_\alpha(M_{x\alpha}\otimes I_R)X_{SR}
(M_{x\alpha}^\dagger\otimes I_R)\succeq0.
\]
这直接证明完全正性。迹的循环性给
$\Tr\mathcal I_x(\rho)=\Tr(\Pi_x\rho)$。
完备关系给 $0\preceq\Pi_x\preceq I$ 与 $\sum_x\Pi_x=I$，因此分支迹不增且总迹保持。
例如两个仪器可取 $M_x=\sqrt{\Pi_x}$ 或 $M_x'=U_x\sqrt{\Pi_x}$，
它们的效应相同，但结果相关的幺正 $U_x$ 会改变条件态。
所以能谱概率不能独自确定后续测量的统计。
\end{proof}

